% Generated draft: requires mathematical and source review.
\documentclass[11pt]{article}
\usepackage[margin=25mm]{geometry}
\usepackage{fontspec}
\setmainfont{Latin Modern Roman}
\usepackage{amsmath,amssymb,mathtools,unicode-math}
\setmathfont{Latin Modern Math}
\usepackage{enumitem,needspace,xurl,hyperref}
\setlength{\parindent}{0pt}
\setlength{\parskip}{5pt}
\setlength{\emergencystretch}{3em}
\newcommand{\N}{\mathbb{N}}
\newcommand{\Z}{\mathbb{Z}}
\newcommand{\Q}{\mathbb{Q}}
\newcommand{\R}{\mathbb{R}}
\newcommand{\C}{\mathbb{C}}
\newcommand{\F}{\mathbb{F}}
\newcommand{\A}{\mathbb{A}}
\newcommand{\PP}{\mathbb{P}}
\newcommand{\E}{\mathbb{E}}
\newcommand{\eps}{\varepsilon}
\newcommand{\dd}{\,\mathrm{d}}
\begin{document}
\section*{UnsolvedMath ID: 6800023}
\textbf{Source problem code:} catalogue entry
\subsection{Definitions and mathematical statement}
Let $V$ be a finite-dimensional Euclidean real vector space, let $\R[V]$ be its polynomial algebra, and put $r^2(v)=\langle v,v\rangle$. In orthonormal coordinates, the Laplacian is $\Delta=\sum_i\partial_i^2$. Under the standard terminology relevant to the cited result, a Laplacian algebra is a unital real subalgebra $A\subseteq\R[V]$ such that
\[
 r^2\in A,\qquad \Delta A\subseteq A.
\]
For $x,y\in V$, define $x\sim_A y$ when $p(x)=p(y)$ for every $p\in A$.

Maximality here means maximality with respect to the induced fibers: every polynomial constant on all $\sim_A$ classes already belongs to $A$. Thus the assertion is
\[
 A=\{q\in\R[V]:
        \forall x,y\in V,\ x\sim_Ay\Longrightarrow q(x)=q(y)\}
\]
for every Laplacian algebra $A$. Equivalently, no strictly larger polynomial subalgebra has exactly the same fibers. This is not the unrelated notion of a maximal proper subalgebra of $\R[V]$.

The short source question supplies no definitions of Laplacian algebra or maximal. The above standard fiber-based definitions should be matched to the cited theorem during review. No homogeneity, connected-fiber, or group-invariant-algebra hypothesis is added.
\subsection{Short English statement}
Must a polynomial algebra containing the squared norm and closed under the Laplacian contain every polynomial constant on its common fibers? The source's intended maximality is fiber-based.

The supplied research summary reports an affirmative solution by R. Mendes and M. Radeschi and quotes Theorem A: every Laplacian algebra $A\subset\R[V]$ is maximal. It identifies this as resolving the earlier question. The registry solved label is retained, while provenance for the 2018 question still marks current status for review. No independent article check is claimed.
\subsection{Sources}
\begin{itemize}
\item Local UnsolvedMath registry, permanent ID 6800023, source code catalogue entry, ``Spherical submetries,'' proposed by Marco Radeschi, Notre Dame, 2018. Dataset curation and metadata: Ulam AI and UnsolvedMath Contributors, CC BY 4.0.
\item Morgan and Pansu, \emph{A List of Open Problems in Differential Geometry} (2018), Question environment 23, \url{https://www.imo.universite-paris-saclay.fr/~pansu/problems_MTDG.tex}; source-TeX extraction, accessed 2026-07-29.
\item R. Mendes and M. Radeschi, ``Maximality of Laplacian algebras, with applications to Invariant Theory,'' \emph{Annali di Matematica Pura ed Applicata} (1923--), 2022, DOI 10.1007/s10231-022-01269-9, Theorem A, as supplied in the research summary. That summary associates arXiv:1903.01532 with its citation; the identifier-to-paper correspondence is not independently verified here.
\item Mendes and Radeschi, ``Laplacian algebras, manifold submetries and the Inverse Invariant Theory Problem,'' \emph{GAFA} (2020), DOI 10.1007/s00039-020-00532-6, as supplied for the earlier conjecture.
\item Source-specific provenance: the public source TeX is on a coauthor's publication page; redistribution terms NEEDS\_REVIEW, separately from the dataset licence.
\end{itemize}
\end{document}
